% =====================  SECTION III  =====================
\section{Proposed Methodology}\label{sec:method}

\subsection{System Overview}\label{subsec:overview}
The system is a linear pipeline of nine stages that carries a user's raw financial
inputs through to an explained, risk-aligned, and audited recommendation. Five raw
inputs are engineered into three interpretable ratios; the ratios are clustered to
produce proxy risk labels; a classifier learns to reproduce those labels from the raw
inputs alone under a feature contract; a deterministic guardrail and a reconciliation
rule enforce conservatism; an occlusion explainer attributes the decision to its
inputs; the reconciled band is mapped to a cryptocurrency allocation, projected over
twelve months, and hashed into an audit record. The complete flow is shown in
Fig.~\ref{fig:arch}. Each stage below is presented with the expression it implements,
its variables, and the property it is designed to guarantee; only methods that the
system actually implements are formalized.

\begin{figure}[t]
\centering
\includegraphics[width=\linewidth]{fig_architecture_v2.png}
\caption{End-to-end architecture, separated into the offline training pipeline (top)
and the online inference pipeline (bottom). During training the five raw inputs are
engineered into three ratios that are standardized and clustered ($k = 3$) to create
the proxy risk labels, whereas the classifier is trained on the five raw inputs only: a
leakage-aware feature contract (Eq.~\ref{eq:contract}) keeps the three ratios out of the
classifier feature matrix, and only the proxy label crosses that boundary, as the
supervised target. The fitted scaler, centroids, preprocessor, and network weights are
then frozen. At inference a user's inputs drive three parallel signals---the frozen MLP
($m$), the frozen cluster assignment ($k$), and the deterministic guardrail
($g$)---which are reconciled to the most cautious band $y^{*} = \max(m, g, k)$
(Eq.~\ref{eq:reconcile}), then explained by occlusion, mapped to a BTC/ETH/USDC
allocation, simulated over twelve months, and written to a hashed, verifiable ledger.}
\label{fig:arch}
\end{figure}

\subsection{Dataset and Provenance}\label{subsec:dataset}
The system is developed and evaluated on a synthetic personal-finance dataset of 3,000
records. The data are generated rather than collected, which permits controlled coverage
of financial scenarios but carries the external-validity caveat discussed in
Section~\ref{sec:discussion}. Two properties of the data are material to the evaluation
and are therefore stated plainly. First, although the file contains 3,000 rows, it holds
only 944 unique user identifiers, so roughly 95.8\% of rows are repeat observations of a
user already present elsewhere in the data; a naive row-level split therefore allows the
same individual to appear in both training and test partitions. Second, each record is
tagged with one of three macro scenarios---normal, inflation, or recession---which shapes
the distribution of incomes and expenses. Table~\ref{tab:dataset} summarizes the
composition and the train/test protocol.

\begin{table}[t]
\centering
\caption{Dataset composition and evaluation protocol}
\label{tab:dataset}
\begin{tabularx}{\linewidth}{@{}l L@{}}
\toprule
Property & Value \\
\midrule
Records / unique users & 3,000 / 944 ($\approx$95.8\% repeat rows) \\
Origin & Synthetic; 28 raw columns \\
Train / test split & 2,400 / 600 (80/20, seed 42) \\
Classifier inputs & 5 raw, user-facing fields \\
Cluster features & 3 engineered ratios \\
Proxy-label balance & Low 44.9\% $\cdot$ Medium 45.9\% $\cdot$ High 9.2\% \\
Scenario tag & normal / inflation / recession \\
Toolchain & scikit-learn 1.5.0, fixed seed \\
\bottomrule
\end{tabularx}
\end{table}

\begin{figure}[t]
\centering
\includegraphics[width=0.55\linewidth]{fig_corr_heatmap.png}
\caption{Pearson correlation among the engineered financial quantities on the synthetic
data. Surplus and expense ratio are strongly anti-correlated ($-0.89$) and income and
surplus are strongly correlated (0.81), reflecting their shared algebraic origin.}
\label{fig:corr}
\end{figure}

\subsection{Feature Engineering}\label{subsec:feature-eng}
Each user supplies monthly income, a credit score, rent or mortgage, essential spending,
discretionary spending, and a loan payment. Total non-debt expenditure is the sum of the
three living-cost components, with the loan payment held out so that debt burden is not
double-counted:
\begin{equation}\label{eq:etotal}
E_{\text{total}} = R + S_{\text{ess}} + S_{\text{disc}}
\end{equation}
Monthly surplus is the income remaining after non-debt expenditure, and is the single
strongest indicator of financial headroom in the system:
\begin{equation}\label{eq:surplus}
\text{surplus} = I - E_{\text{total}}
\end{equation}
The expense ratio expresses non-debt cost as a percentage of income, so that a value at
or above 100\% signals that ordinary spending alone consumes the entire income:
\begin{equation}\label{eq:expratio}
\text{expense\_ratio} = (E_{\text{total}} / I) \times 100
\end{equation}
Debt pressure isolates the loan burden as a percentage of income:
\begin{equation}\label{eq:debt}
\text{debt\_pressure} = (L / I) \times 100
\end{equation}
In Equations~\eqref{eq:etotal}--\eqref{eq:debt}, $I$ is monthly income, $R$ rent or
mortgage, $S_{\text{ess}}$ and $S_{\text{disc}}$ essential and discretionary spending,
and $L$ the loan payment. When income is zero the two ratios are undefined, so the
implementation substitutes a fixed sentinel (999.0) that flags the record as an
extreme-risk outlier rather than dividing by zero. Table~\ref{tab:contract} lists every
variable together with the feature set it belongs to.

\subsection{The Leakage-Aware Feature Contract}\label{subsec:contract}
Because the three ratios are used to construct the labels, handing them to the classifier
would let it trivially recover its own target. The system therefore enforces a feature
contract that partitions the columns into a clustering set and a disjoint classifier set:
\begin{equation}\label{eq:contract}
F_{\text{clu}} \cap F_{\text{cls}} = \emptyset
\end{equation}
with $F_{\text{clu}}$ = \{surplus, expense\_ratio, debt\_pressure\} and $F_{\text{cls}}$
= \{monthly\_income, credit\_score, rent\_or\_mortgage, essential\_spending,
discretionary\_spending\}. A validation routine aborts training if any blocked column
reaches the classifier matrix. This column-level separation is real but incomplete, and
the system reports the gap rather than concealing it: because surplus and expense\_ratio
are algebraic functions of the classifier inputs
(Equations~\eqref{eq:etotal}--\eqref{eq:expratio}), two of the three clustering features
are exactly reconstructible from $F_{\text{cls}}$. Only debt\_pressure is truly withheld,
since it requires the loan payment, which is not a classifier input and is in fact set to
zero at inference. This residual arithmetic leakage is the principal reason a high
classification accuracy is attainable, and it is quantified in
Section~\ref{sec:discussion}.

\begin{table}[t]
\centering
\caption{Variables and the leakage-aware feature contract}
\label{tab:contract}
\begin{tabularx}{\linewidth}{@{}l L L >{\centering\arraybackslash}p{2.3cm}@{}}
\toprule
Symbol & Meaning & Feature set & Recoverable from $F_{\text{cls}}$ \\
\midrule
$I$ & Monthly income & Classifier & --- \\
credit & Credit score & Classifier & --- \\
$R$ & Rent / mortgage & Classifier & --- \\
$S_{\text{ess}}$ & Essential spending & Classifier & --- \\
$S_{\text{disc}}$ & Discretionary spending & Classifier & --- \\
$L$ & Loan payment & Input (zeroed at inference) & --- \\
surplus & $I - E_{\text{total}}$ & Cluster & Yes \\
expense\_ratio & $(E_{\text{total}} / I) \times 100$ & Cluster & Yes \\
debt\_pressure & $(L / I) \times 100$ & Cluster & No (needs $L$) \\
\bottomrule
\end{tabularx}
\end{table}

\subsection{Proxy Label Generation via Clustering}\label{subsec:clustering}
Because the data record no realized default or loss, risk labels are produced
unsupervised. The three ratios live on very different scales, so each is standardized to
zero mean and unit variance, with the statistics estimated on the training partition
only:
\begin{equation}\label{eq:standardize}
z_j = (x_j - \mu_j) / \sigma_j
\end{equation}
K-means then partitions the standardized vectors into $k$ clusters by minimizing the
within-cluster sum of squared Euclidean distances to the cluster centroids, assigning
each point to its nearest centroid:
\begin{equation}\label{eq:kmeans}
J = \sum_{i=1}^{k} \sum_{x \in C_i} \lVert x - \mu_i \rVert^2
\end{equation}
\begin{equation}\label{eq:assign}
c(x) = \arg\min_k \lVert x - \mu_k \rVert^2
\end{equation}
where $C_i$ is the set of points in cluster $i$, $\mu_i$ its centroid, and $J$ the
objective minimized. The cluster count is fixed at $k = 3$ to match the intended
Low/Medium/High banding. This is a documented trade-off rather than an assumption: a
two-cluster solution scores a higher silhouette (0.440), but two clusters cannot express
a three-band policy, so the three-cluster configuration (silhouette 0.370, adjusted Rand
index $0.868 \pm 0.064$ across resampling) is retained. Fitted clusters are ordered into
risk bands by profile---the highest-surplus cluster becomes Low risk and the
highest-expense-ratio cluster becomes High risk.

\subsection{Risk Classification}\label{subsec:classification}
A multilayer perceptron with two hidden layers of 64 and 32 units maps the classifier
feature vector to a risk band. Each hidden layer applies an affine transform and a
rectified-linear activation, and the output layer produces a distribution over the three
classes by softmax:
\begin{equation}\label{eq:h1}
h^{(1)} = \mathrm{ReLU}(W^{(1)}x + b^{(1)})
\end{equation}
\begin{equation}\label{eq:h2}
h^{(2)} = \mathrm{ReLU}(W^{(2)}h^{(1)} + b^{(2)})
\end{equation}
\begin{equation}\label{eq:softmax}
\hat{y}_c = \mathrm{softmax}(W^{(3)}h^{(2)} + b^{(3)})_c
\end{equation}
The network is trained by minimizing the multiclass cross-entropy between the predicted
distribution and the one-hot proxy label over the $N$ training examples:
\begin{equation}\label{eq:celoss}
L = -\frac{1}{N}\sum_{i=1}^{N}\sum_{c=1}^{3} y_{i,c}\log \hat{y}_{i,c}
\end{equation}
where $W^{(l)}$ and $b^{(l)}$ are the weights and biases of layer $l$, $\hat{y}_c$ the
predicted probability of class $c$, and $y_{i,c}$ the one-hot target. The configuration
of the clustering and classifier is given in Table~\ref{tab:config}.

\begin{table}[t]
\centering
\caption{Clustering and classifier configuration}
\label{tab:config}
\begin{tabularx}{\linewidth}{@{}l L@{}}
\toprule
Component & Setting \\
\midrule
Clustering & K-means, $k = 3$, n\_init = 50, seed 42 \\
Model search & K-means \& GMM ($k = 2$--$6$), DBSCAN---16 configs \\
Classifier & MLP, hidden (64, 32), ReLU, softmax(3) \\
Solver / regularization & Adam; L2, $\alpha = 0.001$ \\
Early stopping & val\_fraction 0.15, patience 20, max\_iter 1000 \\
Scaling & StandardScaler fitted on train only \\
\bottomrule
\end{tabularx}
\end{table}

\subsection{Rule-Based Safety Guardrail}\label{subsec:guardrail}
A statistical prediction alone is not a safe basis for advice, so the classifier output
passes through a deterministic guardrail that encodes conservative cash-flow rules and
can only ever raise concern:
\begin{equation}\label{eq:guardrail}
g = \begin{cases}
2 & \text{if } (\text{surplus} < 0 \ \lor\ \text{expense\_ratio} \geq 90),\\
1 & \text{if } \text{expense\_ratio} \geq 60,\\
0 & \text{otherwise.}
\end{cases}
\end{equation}
where 0, 1, and 2 denote Low, Medium, and High risk. A user whose spending exceeds
income, or whose expense ratio reaches 90\%, is treated as High risk irrespective of the
model. The deployed guardrail routine additionally escalates on debt pressure---to High
when debt\_pressure $\geq$ 35\% and to Medium when debt\_pressure $\geq$ 20\%---but on
this dataset those branches are almost entirely subsumed by the cash-flow conditions:
including them leaves the guardrail-only High share unchanged at 60.6\% and shifts the
Low/Medium split by under 0.2 percentage points, so \eqref{eq:guardrail} states the
cash-flow conditions that drive the analysis in Section~\ref{subsec:guardrail-dom}. A
separate recommender module applies a different medium-risk threshold (75 rather than
60); this inconsistency is disclosed as a limitation in Section~\ref{sec:discussion}.

\subsection{Most-Cautious Reconciliation}\label{subsec:reconciliation}
The final band reconciles the model prediction $m$, the guardrail $g$, and the cluster
assignment $k$ by taking the most cautious of the three under the order Low $<$ Medium
$<$ High:
\begin{equation}\label{eq:reconcile}
y^{*} = \max(m, g, k)
\end{equation}
This guarantees that the assessed risk is never lower than any individual signal implies,
so neither a classifier error nor an optimistic cluster can yield an unsafely low
recommendation. Section~\ref{sec:results} quantifies how often each signal is decisive;
the guardrail is the binding signal for most users.

\subsection{Model-Adaptive Explanation}\label{subsec:explanation}
To explain a decision, the system measures how sensitive the predicted probability is to
each input. For the deployed MLP, which exposes no linear coefficients, the influence of
feature $j$ is estimated by occlusion: the feature is replaced by its mean and the
resulting drop in the predicted probability of the assigned class $\hat{c}$ is recorded:
\begin{equation}\label{eq:occlusion}
\Delta_j = \hat{y}_{\hat{c}}(x) - \hat{y}_{\hat{c}}(x_{-j})
\end{equation}
A large influence means the feature strongly supported the prediction. The strategy is
model-adaptive---a linear model would instead report signed coefficients and a tree
ensemble impurity importances---but for the deployed network the mean-imputation
occlusion branch is the one actually exercised. This is an occlusion-based sensitivity
method; it is neither SHAP nor a partial-derivative sensitivity method, and is not
described as such.

\subsection{Risk-Aligned Allocation}\label{subsec:allocation}
The reconciled band maps deterministically to a fixed Bitcoin/Ethereum/USDC allocation
that scales crypto exposure down and stablecoin weighting up as risk rises:
\begin{equation}\label{eq:alloc}
w(\text{Low}) = (35, 35, 30);\quad w(\text{Med}) = (20, 20, 60);\quad w(\text{High}) =
(0, 0, 100)
\end{equation}
with weights in percent summing to 100. A High-risk user is placed entirely in the USDC
stablecoin, a choice motivated by the DeFi-instability literature reviewed in
Section~\ref{sec:related}.

\subsection{Deterministic Twelve-Month Simulation}\label{subsec:simulation}
The allocation is evaluated by compounding an initial capital over twelve months using a
fixed monthly return series for each asset. The portfolio value evolves as
\begin{equation}\label{eq:portfolio}
V_t = V_{t-1}\cdot\left(1 + \sum_a w_a\, r_{a,t}\right)
\end{equation}
and the return over the horizon is the total growth relative to the initial capital:
\begin{equation}\label{eq:roi}
\mathrm{ROI} = (V_{12} - V_0)/V_0
\end{equation}
where $V_t$ is the value at month $t$ ($V_0$ the initial capital), $w_a$ the weight of
asset $a$, and $r_{a,t}$ its fixed monthly return. Because the returns are hardcoded
constants rather than sampled or historical series, this is a deterministic scenario
projection, not a market backtest. The originally reported ROI figures were affected by
an off-by-one normalization (dividing by the first simulated month rather than the
initial capital); Section~\ref{sec:results} reports the corrected values.

\subsection{Tamper-Evident Audit Logging}\label{subsec:audit-log}
Every recommendation is recorded in a tamper-evident ledger. The payload is serialized
into a canonical form with sorted keys and hashed with SHA-256:
\begin{equation}\label{eq:hash}
H = \text{SHA-256}(\,\text{serialize}(\text{payload}, \text{sorted})\,)
\end{equation}
The digest is stored with a unique transaction identifier; integrity verification
re-serializes and re-hashes the stored payload and compares it against $H$, so any
alteration is detectable. The ledger labels entries as a simulated testnet and runs in
memory---it is not deployed to a live blockchain, a limitation revisited in
Section~\ref{sec:discussion}.
